
Quality filters used in large language model pretraining pipelines are designed to select documents matching a notion of textual quality. This paper investigates whether fastText-based quality filtering -- the basis of the DCLM-BASELINE recipe -- also systematically alters demographic-occupation association structure in training corpora. Across seven corpus configurations derived from DCLM-POOL at approximately 50,000-document scale, increasing the fastText percentile threshold from the 10th to the 90th percentile reduced conditional demographic-occupation association entropy H(occupation|demographic) by 22.41\% (Spearman rho = -1.0, p = 1.4 x $10^-24$; bootstrap 95\% CI for H(C5) - H(C1): [-1.154, -0.330]). Conditional log-odds of demographic-occupation co-occurrences across 1,800 pairs increased monotonically with filtering intensity (Spearman rho = 1.0, p = 1.4 x $10^-24$). A proxy-scale model training pilot (H-M2) did not yield a statistically significant graded correlation between corpus entropy and model logit margins (rho = 0.357, p = 0.432), though a negative control comparison showed detectable differences (|C7 - C0| = 0.495). The downstream fairness benchmark evaluation (H-M3) was not executed. These results characterize fastText filtering as a mechanism that restructures demographic associations at the corpus level, though evidence for propagation to model behavior remains inconclusive at the tested scale.

